\subsection{Analyse Statique}\label{sub:Analyse Statique} % (fold)

\begin{definition}{Analyse Statique de Programme}
L'analyse statique de programme correspond à la détermination automatique des propriétés propres à la dynamique d'exécution des programmes dans un contexte statique.
\end{definition}

\paragraph{Fondements}\label{par:Fondements} % (fold)
Au-delà de ce prémisse contextel, le principe de l'analyse est de calculer une sémantique approchée du programme afin de vérifier des spécifications données.
En clair, celà revient à dériver une \emph{sémantique abstraite} calculable et approchée d'une sémantique standard \cite{cousot2001abstract}.
Pour ce faire, on utilise des méthodes formelles, l'interprétaion abstraite, la théorie de la logique, etc.
% paragraph Fondements (end)

%\paragraph{Champ d'action}\label{par:Champ d'action} % (fold)

% paragraph Champ d'action (end)

%\paragraph{Industrialisation et pratiques}\label{par:Industrialisation et pratiques} % (fold)

% paragraph Industrialisation et pratiques (end)

% subsection Analyse Statique (end)
